\section{Related Work}
\label{sec:related}

\textbf{AI dependability and fault taxonomies.}
Avizienis et al.~\cite{avizienis2004} define the foundational taxonomy of
faults, errors, and failures, but it targets conventional deterministic
software: it has no concept of epistemic uncertainty, knowledge coverage,
or reasoning-chain validity, and assumes a complete specification against
which correctness can be verified---an assumption that does not hold for
LLM agents operating with incomplete, potentially stale knowledge.

Humbatova et al.~\cite{humbatova2020} taxonomise real-world bugs in deep
learning systems across model, data, and training faults, but neither the
security domain nor the agentic tool-use paradigm is addressed. Their
taxonomy captures \textit{implementation} faults; \ARES{} captures
\textit{operational} mismatches where the agent is correctly implemented
but insufficiently equipped for the task.

\textbf{LLM hallucination and confidence calibration.}
Huang et al.~\cite{huang2023hallucination} survey hallucination in LLMs but
treat it as a property of the model in isolation, not of the
\textit{operational context} (tool availability, task complexity, knowledge
freshness) that characterises SOC deployments; nor does this literature
examine the multi-agent setting where one agent's hallucination propagates
to downstream reasoning stages.

Guo et al.~\cite{guo2017calibration} show modern neural networks are
systematically overconfident on out-of-distribution inputs. Existing
calibration work measures miscalibration globally; our Epistemic Gap
(Section~\ref{sec:epistemic}) instead targets the specific case where
miscalibration coincides with the point of greatest knowledge weakness.

\textbf{Agent benchmarking and security AI.}
AgentBench~\cite{liu2023agentbench} is the closest existing work to our
evaluation approach, providing a multi-environment benchmark for LLM agents
across web navigation, code execution, and game playing tasks. However,
AgentBench includes no security environment, does not model capability
constraints, and does not support systematic mismatch injection.
PentestGPT~\cite{deng2023pentest} evaluates LLM penetration testing
capabilities, but does not model the conditions under which agents
produce misleading outputs.

The broader AI in cybersecurity literature~\cite{apruzzese2023} catalogues
AI applications across the security lifecycle but explicitly identifies
agent reliability as an open problem. All such work characterises
\textit{what agents can do}, not \textit{how and why they fail}.

\textbf{Contemporary agent-failure taxonomies and benchmarks.}
Two 2024--2025 efforts overlap with \ARES{}'s scope and require
explicit positioning. MAST~\cite{cemri2025mast} derives 14
failure modes from $1{,}600{+}$ annotated traces across seven
multi-agent frameworks and clusters them into system-design,
inter-agent-misalignment, and task-verification categories.
MAST is domain-agnostic and descriptive; \ARES{} is security-specific
and predictive, providing a \textit{pre-execution} reliability
score (ARS) and a formal cascade bound that MAST documents
empirically as ``cascade amplification'' without formalisation.
AgentDojo~\cite{debenedetti2024agentdojo} provides a dynamic
environment with 629 prompt-injection test cases for adversarial
evaluation. \ARES{} covers the orthogonal \textit{non-adversarial}
failure surface: agents failing because capability does not match
task demand, independent of any active attacker. The 2025
hallucination benchmark suite HalluLens~\cite{bang2025hallulens}
measures LLM factual fabrication on generic QA with retrieval
probes, and a concurrent broad survey of LLM-agent
hallucination~\cite{llmagent_halluc_survey2025} catalogues
taxonomy and mitigation approaches across $70{+}$ studies. Both
establish that hallucination is a cross-domain agent failure
mode; \ARES{} contributes the piece that is missing from
this literature: connecting hallucination rate to the specific
capability dimensions of a security task, and modelling cascade
amplification analytically rather than descriptively. At the
application-governance layer, the OWASP Top~10 for Agentic
Applications~\cite{owasp2026agentic} names capability
over-provisioning, tool misuse, and confident hallucination as
top-tier risks; these correspond precisely to $\delta_C$
over-provision, $\delta_C$ and $\delta_R$ interaction, and
Epistemic-Gap failures in our framework.

\textbf{Gap summary.}
Table~\ref{tab:related_isaia} summarises the coverage landscape. To the
best of our knowledge, three gaps are simultaneously unaddressed:
(i)~no formal model covers the full SKC failure surface including
temporal drift and epistemic gap; (ii)~no framework connects capability
mismatch to failure probability quantitatively; and (iii)~no model
predicts hallucination from the interaction of knowledge deficit and
confidence miscalibration. \ARES{} addresses all three.

\begin{table}[t]
\caption{Related Work Coverage Comparison}
\label{tab:related_isaia}
\centering
\footnotesize
\setlength{\tabcolsep}{3pt}
\begin{tabular}{lccccc}
\toprule
\textbf{Work} & \textbf{Formal} & \textbf{Security} & \textbf{Failure} & \textbf{Epistemic} & \textbf{Cascade} \\
& \textbf{Model} & \textbf{Domain} & \textbf{Taxonomy} & \textbf{Gap} & \textbf{Model} \\
\midrule
Avizienis~\cite{avizienis2004}   & \checkmark & -- & \checkmark & -- & -- \\
Huang~\cite{huang2023hallucination} & -- & -- & Partial & -- & -- \\
Humbatova~\cite{humbatova2020}   & -- & -- & \checkmark & -- & -- \\
AgentBench~\cite{liu2023agentbench} & -- & -- & -- & -- & -- \\
PentestGPT~\cite{deng2023pentest} & -- & \checkmark & -- & -- & -- \\
MAST~\cite{cemri2025mast}            & -- & -- & \checkmark & -- & Partial \\
AgentDojo~\cite{debenedetti2024agentdojo} & -- & -- & -- & -- & -- \\
\textbf{\ARES{}} & \checkmark & \checkmark & \checkmark & \checkmark & \checkmark \\
\bottomrule
\end{tabular}
\end{table}
